% !TeX root = ../main.tex
\chapter{TRIỂN KHAI VÀ THỰC NGHIỆM}\label{chap:evaluation}

\section{Môi trường triển khai và tổ chức mã nguồn}
Trên cơ sở kiến trúc phân tách trách nhiệm đã thiết kế ở Chương~\ref{chap:design}, hệ thống Hexta được hiện thực hóa thành tập dịch vụ tái lập đồng nhất giữa môi trường phát
triển cục bộ, cụm fixture kiểm thử cô lập và cấu hình triển khai máy chủ đơn. Bảng~\ref{tab:environment} tổng hợp thông số môi trường và các phiên bản công nghệ chủ chốt trích
xuất trực tiếp từ các tệp khai báo phụ thuộc và cấu hình điều phối trong repository.

\begin{table}[htbp]
  \centering
  \small
  \setlength{\tabcolsep}{4pt}
  \renewcommand{\arraystretch}{1.2}
  \caption{Thông tin môi trường triển khai và ngăn xếp công nghệ}\label{tab:environment}
  \begin{tabularx}{\linewidth}{@{}
    >{\raggedright\arraybackslash}p{3.9cm}
    >{\raggedright\arraybackslash}X@{}}
    \toprule
    Thành phần & Cấu hình và phiên bản xác nhận \\
    \midrule
    Môi trường phát triển và điều phối container &
    Hệ điều hành macOS / Linux; Docker Compose điều phối stack cục bộ qua \texttt{scripts/local.sh} (\texttt{infrastructure/compose.demo.yml} gắn cổng trên
    \texttt{127.0.0.1}),
    cụm fixture cô lập \texttt{infrastructure/compose.connector-test.yml} (\texttt{Postgres :55439}, \texttt{Kafka :59092}), và cấu hình máy chủ đơn
    \texttt{infrastructure/compose.prod.yml} (\texttt{scripts/prod.sh}) định cỡ cho phần cứng 2 vCPU, 16 GB RAM đứng sau \texttt{caddy:2.8-alpine} tự động cấp TLS (chỉ mở cổng 80/443). \\
    Backend và dịch vụ kết nối kênh (Go) &
    \texttt{go.work} dùng Go 1.26.2 cho 4 module (\texttt{packages/shared}, \texttt{services/api}, \texttt{services/connector}, \texttt{test}). \texttt{services/api} sử dụng Gin v1.12.0,
    Uber Fx v1.24.0, Uber Zap v1.27.1, GORM v1.31.1, \texttt{grpc} v1.79.1, \texttt{go-redis/v9} v9.18.0, \texttt{golang-jwt/v5} v5.2.0, \texttt{shopspring/decimal} v1.4.0 và
    Swaggo \texttt{swag} v1.16.6;
    \texttt{services/connector} sử dụng \texttt{pgx/v5} v5.6.0 và \texttt{segmentio/kafka-go} v0.4.49. \\
    Dịch vụ thực thi AI (Python) &
    Python 3.12--3.14 (image \texttt{python:3.14-slim}, trình quản lý \texttt{uv} 0.11.11); FastAPI 0.143.5, Uvicorn 0.54.0, \texttt{grpcio} 1.84.0, Protobuf $\ge$7.36.2,
    LangGraph 1.2.16,
    LangChain 1.4.8 tích hợp Gemini (\texttt{langchain-google-genai}), OpenAI (\texttt{langchain-openai}), Anthropic (\texttt{langchain-anthropic}) và chế độ \texttt{fake};
    \texttt{psycopg\_pool} 3.3.3; kiểm thử bằng \texttt{pytest}. \\
    Giao diện người dùng và SDK &
    Next.js 16.3.0, React 19.2.8, TypeScript 5, Tailwind CSS 4, Zustand 5.0.15, \texttt{react-markdown}, quản lý workspace bằng \texttt{pnpm} 10.x;
    gói \texttt{@hexta/sdk} dùng Axios 1.19.0, \texttt{tsup} 8.5.1, Vitest 5.0.0 và thư viện giao diện \texttt{@hexta/ui}. \\
    Cơ sở dữ liệu, bộ nhớ đệm và thông điệp &
    PostgreSQL \texttt{17.7-alpine3.23}, Redis \texttt{7.4-alpine}, Kafka \texttt{cp-\hspace{0pt}kafka:\hspace{0pt}7.6.1} của Confluent (KRaft một nút, ghi đĩa \texttt{fsync} từng thông điệp, tắt tự tạo topic), và công cụ quản lý
    lược đồ Atlas CLI (\texttt{arigaio/atlas}). \\
    \bottomrule
  \end{tabularx}
\end{table}

Trong cấu hình cục bộ (\texttt{compose.\hspace{0pt}demo.\hspace{0pt}yml}), các dịch vụ chỉ lắng nghe trên \texttt{127.0.0.1} (\texttt{web} cổng \texttt{3000}, \texttt{api} cổng \texttt{8080},
\texttt{ai} HTTP \texttt{58000} và gRPC \texttt{55051}, \texttt{postgres} \texttt{55433}, \texttt{redis} \texttt{56379}), đồng thời AI mặc định dùng chế độ \texttt{fake} để
tái lập không cần khóa trả phí.
Ở cấu hình máy chủ đơn (\texttt{compose.\hspace{0pt}prod.\hspace{0pt}yml} và \texttt{infrastructure/\hspace{0pt}prod/\hspace{0pt}Caddyfile}), chỉ \texttt{caddy} phơi cổng công khai \texttt{80/443}: mọi đường dẫn
\texttt{/api/*} chuyển tới \texttt{api:8080}, chỉ hai điểm nhận \texttt{/webhooks/\hspace{0pt}facebook} và \texttt{/webhooks/\hspace{0pt}telegram/*} (giới hạn thân 1 MB) định tuyến tới
\texttt{connector:8091}, các đường dẫn \texttt{/webhooks/*} khác trả về 404, phần còn lại chuyển tới \texttt{web:3000}.
Mặt phẳng điều khiển \texttt{/internal/*} của Connector, dịch vụ AI, PostgreSQL, Redis và Kafka hoàn toàn ẩn trong mạng nội bộ Compose.
Hệ thống khởi tạo tường minh theo cơ chế lũy đẳng năm topic Kafka (3 phân vùng mỗi topic) gồm: \texttt{hexta.\hspace{0pt}ingress.\hspace{0pt}received.\hspace{0pt}v1}, \texttt{hexta.\hspace{0pt}outbound.\hspace{0pt}intent.\hspace{0pt}v1}, \texttt{hexta.\hspace{0pt}outbound.\hspace{0pt}result.\hspace{0pt}v1}, \texttt{hexta.\hspace{0pt}connector.\hspace{0pt}health.\hspace{0pt}v1} và \texttt{hexta.\hspace{0pt}orders.\hspace{0pt}events.\hspace{0pt}v1}.

Mã nguồn toàn hệ thống được tổ chức theo mô hình monorepo gồm các phân vùng sở hữu rõ ràng:
\begin{itemize}
  \item \texttt{apps/web}: Ứng dụng giao diện Next.js (App Router) phục vụ vận hành doanh nghiệp, bán hàng, quy trình, hộp thư hội thoại, cấu hình kênh và tư vấn mua sắm (cùng
    các ứng dụng định hướng mở rộng \texttt{admin}, \texttt{docs}, \texttt{landing}).
  \item \texttt{services/api}: Dịch vụ nghiệp vụ trung tâm bằng Go với các điểm vào \texttt{cmd/\hspace{0pt}main.go} (HTTP API), \texttt{cmd/\hspace{0pt}ingress} (tiến trình tiêu thụ Kafka độc lập),
    \texttt{cmd/\hspace{0pt}platform-\hspace{0pt}operator} (cấp quyền vận hành nền tảng) và \texttt{cmd/\hspace{0pt}ai-\hspace{0pt}smoke}.
  \item \texttt{services/\hspace{0pt}connector}: Dịch vụ kết nối kênh giao tiếp ngoài bằng Go (\texttt{cmd/\hspace{0pt}connector}), độc lập về quy trình và lưu trữ bí mật.
  \item \texttt{services/ai}: Dịch vụ thực thi AI bằng Python cung cấp hai điểm cuối HTTP nội bộ \texttt{/internal/\hspace{0pt}v1/\hspace{0pt}preview} và \texttt{/internal/\hspace{0pt}v1/\hspace{0pt}consult} cùng máy chủ gRPC trên cổng \texttt{50051}.
  \item \texttt{packages/\hspace{0pt}shared}, \texttt{packages/\hspace{0pt}sdk}, \texttt{packages/\hspace{0pt}ui}: Các gói dùng chung gồm chuẩn lỗi, hợp đồng sự kiện \texttt{ingress}/\texttt{outbound} v1 và
    \texttt{kafkatransport} cho Go; client \texttt{@hexta/sdk} và bộ thành phần giao diện \texttt{@hexta/ui} cho TypeScript.
  \item \texttt{proto/\hspace{0pt}hexta/\hspace{0pt}ai/\hspace{0pt}v1}, \texttt{migrations/\hspace{0pt}\{api,ai,connector\}}, \texttt{infrastructure/}: Đặc tả Protobuf \texttt{service.proto}, lịch sử migration Atlas cho ba
    phân vùng dữ liệu (17 tệp \texttt{api}, 1 tệp \texttt{ai}, 4 tệp \texttt{connector}), cùng cấu hình Compose và tập lệnh \texttt{scripts/}.
\end{itemize}

Nhằm bảo đảm biên cô lập dữ liệu ngay tại tầng lưu trữ, \texttt{infrastructure/\hspace{0pt}prod/\hspace{0pt}init-\hspace{0pt}db.sh} thiết lập ba tài khoản cơ sở dữ liệu tách biệt trên PostgreSQL:
(1) \texttt{hexta\_\hspace{0pt}api} sở hữu cơ sở dữ liệu \texttt{hexta} (lược đồ \texttt{public} và bảng lịch sử \texttt{atlas\_\hspace{0pt}schema\_\hspace{0pt}revisions});
(2) \texttt{hexta\_\hspace{0pt}ai} chỉ sở hữu lược đồ \texttt{ai\_\hspace{0pt}runtime} (và \texttt{atlas\_\hspace{0pt}ai\_\hspace{0pt}revisions}) bên trong \texttt{hexta} để quản lý tác vụ AI;
và (3) \texttt{hexta\_\hspace{0pt}connector} sở hữu cơ sở dữ liệu độc lập \texttt{hexta\_\hspace{0pt}connector} sau khi thu hồi quyền \texttt{CONNECT} từ \texttt{PUBLIC}, bảo đảm API không thể truy
cập trực tiếp khóa mã hóa kênh và việc chạy Atlas trên từng dịch vụ không xung đột bảng ngoại lai.

\section{Dữ liệu và kịch bản thực nghiệm}
\subsection{Bộ dữ liệu thực nghiệm và mô phỏng đối tác}
Để kiểm chứng hoạt động của hệ thống sát với bối cảnh thực tế mà không ảnh hưởng đến giao dịch thật, bộ dữ liệu thực nghiệm được thiết kế dựa trên nghiệp vụ của Doanh nghiệp A
(kinh doanh cây ăn quả và cây cảnh -- Chương~\ref{chap:introduction}), kết hợp các tình huống biên và máy chủ giả lập được gắn nhãn fixture/mô phỏng tường minh:
\begin{itemize}
  \item \textbf{Dữ liệu đa doanh nghiệp (\texttt{Tenant A} và \texttt{Tenant B}):} \texttt{Tenant A} đại diện cho doanh nghiệp cây ăn quả/cây cảnh với tài khoản chủ sở hữu,
    nhân viên tư vấn bán hàng và nhân viên đóng gói/vận chuyển;
  \texttt{Tenant B} đóng vai trò doanh nghiệp đối chứng để kiểm tra cô lập tài nguyên chéo, thu hồi quyền thành viên và ràng buộc duy nhất một trang Facebook hoặc một bot Telegram chỉ gắn với một doanh nghiệp đang hoạt động.
  \item \textbf{Danh mục sản phẩm và các trường hợp biên về giá:} Danh mục gồm các biến thể đã khai báo đơn giá minh bạch (\texttt{price\_provided = true}, ví dụ các giống cây theo kích thước bầu),
  biến thể có giá bằng không tường minh (\texttt{price\_provided = true}, \texttt{price = 0} cho vật tư/cây tặng kèm) và biến thể chưa có giá (\texttt{price\_provided =
  false}) nhằm kiểm chứng cổng chặn xác nhận đơn và bộ lọc danh mục đưa vào tư vấn viên AI.
  \item \textbf{Đơn nháp và quy trình công việc hỗn hợp:} Tập đơn nháp bao gồm các đơn chưa đủ thông tin, đơn cần giao hàng (\texttt{requires\_delivery = true} bắt buộc có
    \texttt{delivery\_address}) đối chứng với đơn nhận trực tiếp (\texttt{requires\_delivery = false}),
  gắn với quy trình tuần tự nhiều bước pha trộn cả chế độ chuyển trực tiếp (\texttt{direct}) và chờ xác nhận nhận (\texttt{acceptance}).
  \item \textbf{Máy chủ giả lập Meta Graph API và Telegram Bot API:} Hai máy chủ HTTP cục bộ (\texttt{mockGraph} và \texttt{mockTelegram}) dựng trên \texttt{httptest.Server}
    mô phỏng đầy đủ luồng cấp phép, webhook có chữ ký/bí mật và các nhánh lỗi mạng:
  \texttt{mockGraph} trả về thành công hoặc mã lỗi hết cửa sổ 24 giờ (\texttt{code = 10}, \texttt{subcode = 2018278} $\rightarrow$ \texttt{window\_24h}) và thu hồi token (\texttt{code = 190} $\rightarrow$ \texttt{auth\_expired});
  \texttt{mockTelegram} phân nhánh theo \texttt{chat\_id} (\texttt{9001} thành công, \texttt{9002} HTTP 403 \texttt{bot\_blocked}, \texttt{9003} HTTP 400 \texttt{chat\_not\_found}, \texttt{9004} đóng kết nối TCP đột ngột trả về \texttt{unknown}, \texttt{9005} HTTP 401 \texttt{auth\_expired}).
  \item \textbf{Dữ liệu kiểm thử AI có kiểm soát:} Gồm tập kịch bản mẫu đang bật/tắt, facts doanh nghiệp, truy vấn lọc giá tiếng Việt và phản hồi giả lập ở tầng vận chuyển
    HTTP cho Gemini, OpenAI, Anthropic cùng chế độ \texttt{FakeProvider} tất định.
\end{itemize}

\subsection{Thiết kế các nhóm kịch bản thực nghiệm}
Toàn bộ thực nghiệm kiểm chứng được chia thành bốn nhóm kịch bản có kiểm soát (K1--K4), ánh xạ trực tiếp tới ba câu hỏi nghiên cứu (RQ1, RQ2, RQ3) ở
Chương~\ref{chap:introduction} như tóm tắt trong Bảng~\ref{tab:scenarios}.

\begin{table}[htbp]
  \centering
  \small
  \setlength{\tabcolsep}{4pt}
  \renewcommand{\arraystretch}{1.2}
  \caption{Ánh xạ bốn nhóm kịch bản thực nghiệm (K1--K4) tới các câu hỏi nghiên cứu}\label{tab:scenarios}
  \begin{tabularx}{\linewidth}{@{}
    >{\raggedright\arraybackslash}p{0.9cm}
    >{\raggedright\arraybackslash}p{3.3cm}
    >{\raggedright\arraybackslash}X
    >{\raggedright\arraybackslash}p{1.6cm}@{}}
    \toprule
    Mã & Nhóm kịch bản & Trọng tâm và điều kiện biên được kiểm chứng & Câu hỏi \\
    \midrule
    K1 & Onboarding, phân quyền và cô lập đa doanh nghiệp &
    Xét duyệt đơn đăng ký qua \texttt{platform-operator}, mời thành viên băm token, làm mới phiên qua cookie \texttt{HttpOnly}, chặn truy cập chéo \texttt{Tenant
    A}/\texttt{Tenant B} và khóa ngoại tổ hợp. &
    RQ3 \\
    K2 & Đơn nháp, cổng xác nhận và quy trình bàn giao &
    Lưu đơn nháp thiếu trường, phân biệt \texttt{price\_provided} với giá 0, kiểm tra giao hàng, bàn giao hỗn hợp \texttt{direct}/\texttt{acceptance}, bất biến phiên bản và
    tranh chấp \texttt{accept}/\texttt{assign}. &
    RQ1, RQ3 \\
    K3 & Hộp thư đa kênh, điều khiển Desk và đồng bộ Kafka &
    Chuyển trạng thái \texttt{BOT\_ACTIVE} / \texttt{WAITING\_STAFF} / \texttt{STAFF\_ACTIVE} với \texttt{control\_generation}; Facebook OAuth và Telegram Bot; Outbox
    $\rightarrow$ Kafka $\rightarrow$ Inbox; gửi một lần và các lỗi \texttt{window\_24h}, \texttt{bot\_blocked}, \texttt{chat\_not\_found}, \texttt{auth\_expired},
    \texttt{unknown}. &
    RQ2, RQ3 \\
    K4 & Thực thi AI bền vững và tư vấn mua sắm SSE &
    \texttt{RunStore} thu hồi \texttt{lease}, hàng rào \texttt{LLM\_CALL\_STARTED}, hủy \texttt{CancelRun}, kiểm chứng đầu ra LangGraph và luồng tư vấn SSE \texttt{/api/v1/ai/converse} nối sang \texttt{checkout}. &
    RQ2, RQ3 \\
    \bottomrule
  \end{tabularx}
\end{table}

Chi tiết các ca biên và tình huống xung đột được kiểm chứng ở từng nhóm như sau:
\begin{itemize}
  \item \textbf{Nhóm K1:} Kiểm tra từ chối đơn đăng ký thiếu lý do, ngăn người không phải thành viên hoặc thành viên đã bị thu hồi quyền đọc/ghi tài nguyên;
  xác minh phiên đăng nhập truy vấn trực tiếp PostgreSQL thu hồi hiệu lực ngay; kiểm tra khóa ngoại tổ hợp \texttt{(tenant\_\hspace{0pt}id, id)} chặn ghép chéo đơn hàng hoặc biến thể giữa
  \texttt{Tenant A} và \texttt{Tenant B}.
  \item \textbf{Nhóm K2:} Kiểm tra lưu đơn nháp không cần đủ trường, từ chối xác nhận khi \texttt{price\_\hspace{0pt}provided = false} hoặc thiếu \texttt{delivery\_\hspace{0pt}address} khi giao hàng,
    nhưng chấp nhận khi \texttt{price = 0} với \texttt{price\_\hspace{0pt}provided = true};
  kiểm tra xuất bản quy trình mới không làm thay đổi \texttt{work\_\hspace{0pt}items} đã khởi tạo; chạy đồng thời hai luồng sửa/xác nhận đơn và cạnh tranh giữa \texttt{accept} với
  \texttt{assign} để xác nhận tuần tự hóa theo khóa hàng và phiên bản.
  \item \textbf{Nhóm K3:} Kiểm tra chuyển đổi quyền điều khiển Desk, từ chối trả lời khi \texttt{context\_\hspace{0pt}revision} cũ hoặc ở trạng thái \texttt{UNKNOWN};
  kiểm tra vòng đời OAuth Facebook và Bot Token Telegram (1 kết nối = 1 doanh nghiệp, \texttt{rotate\_\hspace{0pt}token} khác bot trả về 409, sai secret trả về 403, bỏ qua tin nhóm,
  chuyển media thành placeholder);
  tắt tiến trình Connector khi Kafka mất kết nối rồi khôi phục để chứng minh truyền đủ sự kiện webhook, khử trùng lặp và chỉ gửi ra nhà cung cấp đúng 1 lần cho mỗi
  \texttt{attempt\_\hspace{0pt}id} kể cả khi gặp lỗi \texttt{unknown}.
  \item \textbf{Nhóm K4:} Kiểm tra xung đột đầu vào \texttt{StartRun} khác digest, gia hạn và thu hồi \texttt{lease} tối đa 3 lần thử (\texttt{RECOVERY\_\hspace{0pt}ATTEMPTS\_\hspace{0pt}EXHAUSTED}),
    chặn gọi lại LLM thật sau dấu \texttt{LLM\_\hspace{0pt}CALL\_\hspace{0pt}STARTED} (\texttt{LLM\_\hspace{0pt}RECOVERY\_\hspace{0pt}UNCERTAIN}), hủy qua \texttt{CancelRun};
  kiểm chứng \texttt{graph.py} chuyển \texttt{human} khi thiếu dữ liệu hoặc bị tắt \texttt{allow\_\hspace{0pt}generation}, và xác minh \texttt{ConsultantService} loại bỏ sản phẩm/biến thể
  giả từ phản hồi AI trước khi phát SSE.
\end{itemize}

\section{Cài đặt các chức năng chính}
\subsection{Phân hệ quản lý doanh nghiệp, bán hàng và quy trình công việc}
Dịch vụ \texttt{services/api} được lắp ghép qua Uber Fx (\texttt{internal/\hspace{0pt}bootstrap}) thành bốn lớp tách biệt: controller HTTP (Gin), service nghiệp vụ, repository (GORM trên
PostgreSQL với bể kết nối tối đa 50 kết nối mở, 25 nhàn rỗi, thời gian sống 30 phút) và hạ tầng ngoài.
Nhóm đường dẫn doanh nghiệp \texttt{/api/\hspace{0pt}v1/\hspace{0pt}tenants/\hspace{0pt}:tenantID/*} áp dụng cặp middleware \texttt{RequireAuth} và \texttt{RequireTenantMember} (ủy quyền
\texttt{TenantService.\hspace{0pt}AuthorizeMember} đọc trực tiếp DB và đối chiếu khớp tiêu đề \texttt{X-\hspace{0pt}Tenant-\hspace{0pt}ID}), đồng thời \texttt{ISessionRepository} gắn trực tiếp với PostgreSQL để
việc thu hồi phiên có hiệu lực tức thì.
\texttt{TenantOnboardingService} tách quyền xét duyệt của \texttt{platform-\hspace{0pt}operator} khỏi chủ doanh nghiệp, còn lời mời thành viên chỉ lưu mã băm SHA-256 của token.

Ở phân hệ bán hàng (\texttt{SalesService}), lệnh \texttt{SaveDraft} cho phép lưu đơn nháp bán phần (tối đa 100 dòng, số lượng dương đến 1.000.000, tiền tệ không âm dưới
$10^{13}$ làm tròn 2 chữ số thập phân) mà không tự điền đơn giá vắng mặt.
Tại lệnh \texttt{Confirm}, giao dịch khóa đơn, kiểm tra đồng thời người phụ trách, trạng thái \texttt{draft}, \texttt{version} lạc quan, thông tin liên hệ khách hàng, cờ
\texttt{price\_\hspace{0pt}provided = true} trên từng dòng hàng, địa chỉ nếu \texttt{requires\_\hspace{0pt}delivery = true}.
Chỉ khi mọi điều kiện hợp lệ, hệ thống mới cập nhật trạng thái \texttt{confirmed}, ghi \texttt{audit\_\hspace{0pt}logs} và chèn bản ghi \texttt{order.\hspace{0pt}confirmed} vào bảng
\texttt{domain\_\hspace{0pt}events} trong cùng giao dịch.
Tính năng tìm kiếm sản phẩm dùng cột sinh \texttt{products.\hspace{0pt}search\_\hspace{0pt}vector} kèm chỉ mục GIN trên PostgreSQL thay cho cụm tìm kiếm ngoài (Elasticsearch/Qdrant chuyển sang
profile tùy chọn \texttt{search-\hspace{0pt}engine}).

Ở phân hệ quy trình (\texttt{WorkflowService} và gói \texttt{repository/\hspace{0pt}workflowmodel}), chủ doanh nghiệp xuất bản chuỗi từ 1 đến 30 bước với chính sách bàn giao
\texttt{direct} hoặc \texttt{acceptance} vào các bảng bất biến \texttt{workflow\_\hspace{0pt}versions} và \texttt{workflow\_\hspace{0pt}steps}.
Khi khởi tạo \texttt{work\_\hspace{0pt}items} từ đơn xác nhận, công việc gắn cố định \texttt{workflow\_\hspace{0pt}version\_\hspace{0pt}id} và \texttt{current\_\hspace{0pt}step}; với bước \texttt{acceptance}, lệnh
\texttt{Handoff} chỉ ghi \texttt{recipient\_\hspace{0pt}id} và giữ nguyên người chịu trách nhiệm hiện tại (\texttt{responsible\_\hspace{0pt}id}) cho đến khi người được chỉ định gọi \texttt{Accept}, trong khi
\texttt{Assign} của chủ doanh nghiệp hủy đề cử chờ.
Mọi chuyển đổi sử dụng khóa hàng \texttt{FOR UPDATE}, kiểm tra \texttt{expected\_\hspace{0pt}version}, chống tạo trùng bằng \texttt{Idempotency-\hspace{0pt}Key} kèm hàm băm SHA-256 và ghi lịch sử bất
biến vào \texttt{work\_\hspace{0pt}events}.

\subsection{Phân hệ hộp thư, điều khiển Desk và đồng bộ sự kiện Kafka}
Phân hệ hội thoại (\texttt{MessagingService}) quản lý \texttt{messaging\_\hspace{0pt}threads}, \texttt{messaging\_\hspace{0pt}turns} (ràng buộc duy nhất một lượt \texttt{open} mỗi hội thoại) và
\texttt{messaging\_\hspace{0pt}messages}.
Hàm thuần \texttt{Handling\hspace{0pt}State()} chiếu trực tiếp các trường bền vững (\texttt{Control}, \texttt{Assignee\hspace{0pt}ID}, \texttt{Needs\hspace{0pt}Human}) thành bốn trạng thái: \texttt{BOT\_\hspace{0pt}ACTIVE},
\texttt{WAITING\_\hspace{0pt}STAFF}, \texttt{STAFF\_\hspace{0pt}ACTIVE} hoặc cơ chế đóng an toàn \texttt{UNKNOWN}.
Các lệnh \texttt{claim}, \texttt{assign}, \texttt{handoff} (yêu cầu mã lý do hợp lệ như \texttt{customer\_\hspace{0pt}requested\_\hspace{0pt}staff}, \texttt{insufficient\_\hspace{0pt}information},
\texttt{ai\_\hspace{0pt}error}, \texttt{business\_\hspace{0pt}review}, \texttt{staff\_\hspace{0pt}handoff}) và lệnh \texttt{return-\hspace{0pt}to-\hspace{0pt}bot} đều khóa \texttt{messaging\_\hspace{0pt}threads} qua \texttt{messagingLockedTurn}, kiểm
tra \texttt{expected\_\hspace{0pt}version}, tăng đồng thời \texttt{Version} và \texttt{Control\hspace{0pt}Generation}, rồi ghi vết trước/sau vào \texttt{audit\_\hspace{0pt}logs}.

Khi nhân viên được phân công trả lời kênh thật qua \texttt{SendChannel} (\texttt{POST} tới \texttt{/api/\hspace{0pt}v1/\hspace{0pt}tenants/\hspace{0pt}\{tenantID\}/\hspace{0pt}turns/\hspace{0pt}\{id\}/\hspace{0pt}messages}), hệ thống kiểm tra độ dài tối đa
2.000 ký tự (\texttt{outbound.\hspace{0pt}MaxTextRunes}), quyền nhân viên phụ trách, trạng thái \texttt{STAFF\_\hspace{0pt}ACTIVE}, kết nối đang \texttt{active}, khớp \texttt{expected\_\hspace{0pt}version} và
\texttt{context\_\hspace{0pt}revision}, cùng khóa chống lặp \texttt{client\_\hspace{0pt}message\_\hspace{0pt}id} kèm SHA-256 hash.
Sau đó, giao dịch ghi nguyên tử tin nhắn (\texttt{status = "channel"}, \texttt{delivery\_\hspace{0pt}status = "sending"}), \texttt{outbound\_\hspace{0pt}intents} (TTL 10 phút),
\texttt{outbound\_\hspace{0pt}attempts} và sự kiện \texttt{outbound.\hspace{0pt}intent} vào \texttt{domain\_\hspace{0pt}events}.

Để tách mạng khỏi giao dịch cơ sở dữ liệu, \texttt{OutboxRelay} thuê khóa sự kiện từ \texttt{domain\_\hspace{0pt}events} bằng \texttt{FOR UPDATE SKIP LOCKED} kèm hàng rào \texttt{fence}
và lease 30 giây (tối đa 12 lần thử với backoff lũy thừa tối đa 300 giây), phát \texttt{order.\hspace{0pt}confirmed} tới \texttt{hexta.\hspace{0pt}orders.\hspace{0pt}events.\hspace{0pt}v1} và \texttt{outbound.\hspace{0pt}intent} tới
\texttt{hexta.\hspace{0pt}outbound.\hspace{0pt}intent.\hspace{0pt}v1} theo khóa phân vùng \texttt{ConnectionID:\hspace{0pt}RecipientID}.
Ở chiều vào, tiến trình \texttt{cmd/\hspace{0pt}ingress} chạy ba consumer group cho \texttt{hexta.\hspace{0pt}ingress.\hspace{0pt}received.\hspace{0pt}v1}, \texttt{hexta.\hspace{0pt}outbound.\hspace{0pt}result.\hspace{0pt}v1} và
\texttt{hexta.\hspace{0pt}connector.\hspace{0pt}health.\hspace{0pt}v1}; offset chỉ được commit sau khi giao dịch PostgreSQL thành công, còn thông điệp sai cấu trúc được cách ly theo giá trị băm vào bảng
\texttt{messaging\_\hspace{0pt}ingress\_\hspace{0pt}quarantine}.

\subsection{Dịch vụ kết nối kênh giao tiếp ngoài}
Dịch vụ \texttt{services/\hspace{0pt}connector} cài đặt bộ hợp đồng plugin (\texttt{Plugin}, \texttt{Webhook}, \texttt{Provisioner}, \texttt{BotProvisioner}, \texttt{TextSender},
\texttt{BasePlugin}) tại \texttt{internal/\hspace{0pt}plugin/\hspace{0pt}plugin.go} cùng lớp lưu trữ chịu tải cao (\texttt{connector.\hspace{0pt}receipts} khử lặp theo biên nhận, giới hạn hàng đợi xấp xỉ 100.000
sự kiện chờ, chèn khóa theo thứ tự sắp xếp để tránh deadlock):
\begin{itemize}
  \item \textbf{Adapter Facebook (\texttt{plugin/\hspace{0pt}facebook}):} Kiểm tra chữ ký HMAC-SHA256 trên raw body (\texttt{X-\hspace{0pt}Hub-\hspace{0pt}Signature-\hspace{0pt}256}), xử lý \texttt{hub.\hspace{0pt}challenge}, trao đổi
    mã OAuth $\rightarrow$ user token dài hạn $\rightarrow$ Page token qua Meta Graph API \texttt{v21.0}, lưu tạm trong \texttt{connector.\hspace{0pt}oauth\_\hspace{0pt}sessions} (TTL 10 phút), đăng
    ký \texttt{subscribed\_\hspace{0pt}apps} và ánh xạ đính kèm không hỗ trợ thành placeholder \texttt{unsupported\_\hspace{0pt}media} mà không tải tệp nhị phân.
  \item \textbf{Adapter Telegram (\texttt{plugin/\hspace{0pt}telegram}):} Xác thực Bot Token bằng \texttt{getMe}, sinh bí mật 32 byte cho đường dẫn webhook riêng
    \texttt{/webhooks/\hspace{0pt}telegram/\hspace{0pt}\{connectionId\}}, so sánh tiêu đề \texttt{X-\hspace{0pt}Telegram-\hspace{0pt}Bot-\hspace{0pt}Api-\hspace{0pt}Secret-\hspace{0pt}Token} theo thời gian hằng số (\texttt{subtle.\hspace{0pt}ConstantTimeCompare}), lọc bỏ
    các cập nhật không phải chat \texttt{private} trước khi đưa vào outbox và gọi \texttt{deleteWebhook} khi tạm dừng hoặc ngắt kết nối.
  \item \textbf{Niêm phong AES-256-GCM AAD và gửi một lần:} \texttt{secret.\hspace{0pt}Sealer} mã hóa toàn bộ token và webhook secret bằng khóa 32 byte
    (\texttt{CONNECTOR\_\hspace{0pt}CREDENTIAL\_\hspace{0pt}KEY\_\hspace{0pt}FILE}) ràng buộc dữ liệu liên kết như \texttt{oauth\_\hspace{0pt}session:\hspace{0pt}\{id\}}, \texttt{connection:\hspace{0pt}\{id\}} và \texttt{webhook\_\hspace{0pt}secret:\hspace{0pt}\{id\}}.
  Khi gửi (\texttt{Outbound\hspace{0pt}Service.\hspace{0pt}Handle}), \texttt{Begin\hspace{0pt}Attempt} chèn trạng thái \texttt{dispatching}; mọi lần cấp lại trên trạng thái này đều chuyển sang \texttt{unknown}
  (\texttt{interrupted}) và tuyệt đối không gọi lại API ngoài.
  Lỗi trả về được chuẩn hóa sạch bí mật thành \texttt{window\_\hspace{0pt}24h}, \texttt{bot\_\hspace{0pt}blocked}, \texttt{chat\_\hspace{0pt}not\_\hspace{0pt}found}, \texttt{auth\_\hspace{0pt}expired}, \texttt{rate\_\hspace{0pt}limit},
  \texttt{invalid\_\hspace{0pt}recipient} hoặc \texttt{transport} (gắn cờ \texttt{Unknown} khi lỗi mạng hoặc HTTP $\ge 500$).
\end{itemize}

\subsection{Dịch vụ AI bền vững, giao diện người dùng và SDK}
Dịch vụ \texttt{services/ai} chạy FastAPI trên Uvicorn song song với máy chủ gRPC \texttt{AIService} (\texttt{StartRun}, \texttt{GetRun}, \texttt{CancelRun}) dùng chung
\texttt{RunStore} trên \texttt{psycopg\_\hspace{0pt}pool} (2--10 kết nối, timeout checkout 5 giây).
Yêu cầu \texttt{StartRun} được băm tất định (SHA-256 trên Protobuf) để phát hiện xung đột đầu vào trên cùng \texttt{run\_\hspace{0pt}id}.
Worker nền dùng \texttt{FOR UPDATE SKIP LOCKED} thuê tác vụ (mặc định 60 giây, tối đa 3 lần thử); trước khi gọi LLM trả phí, \texttt{begin\_\hspace{0pt}paid\_\hspace{0pt}execution} ghi dấu
\texttt{LLM\_\hspace{0pt}CALL\_\hspace{0pt}STARTED} có điều kiện \texttt{lease\_\hspace{0pt}token} nhằm ngăn lặp cuộc gọi nếu phục hồi sự cố (\texttt{LLM\_\hspace{0pt}RECOVERY\_\hspace{0pt}UNCERTAIN}).
\texttt{LangChainProvider} hợp nhất Gemini, OpenAI và Anthropic qua \texttt{init\_\hspace{0pt}chat\_\hspace{0pt}model} và \texttt{with\_\hspace{0pt}structured\_\hspace{0pt}output}, tắt LangSmith tracing (\texttt{LANGSMITH\_\hspace{0pt}TRACING=false}) và tắt tự động thử lại ở SDK.

Hai đồ thị LangGraph gồm quy trình định tuyến ứng viên (\texttt{graph.py}: \texttt{validate\_\hspace{0pt}input} $\rightarrow$ \texttt{prepare\_\hspace{0pt}context} $\rightarrow$
\texttt{classify\_\hspace{0pt}intent} $\rightarrow$ các nhánh \texttt{script}, \texttt{generate}, \texttt{clarify} hoặc \texttt{human} $\rightarrow$ \texttt{validate\_\hspace{0pt}output}) và quy trình tư vấn mua sắm
(\texttt{graph\_\hspace{0pt}consultant.py} áp dụng guardrail đầu vào, hàm trích xuất khoảng giá tiếng Việt \texttt{parse\_\hspace{0pt}vietnamese\_\hspace{0pt}price\_\hspace{0pt}filter} và ràng buộc chỉ tư vấn từ danh mục
được cấp).
Ở phía Go, \texttt{ConsultantService} nạp tối đa 100 sản phẩm / 200 biến thể có \texttt{price\_\hspace{0pt}provided = true}, gọi \texttt{/internal/\hspace{0pt}v1/\hspace{0pt}consult}, tái xác thực quyền thành
viên, biến thể, đơn giá gốc và liên kết ảnh \texttt{https} hợp lệ trước khi phát khung SSE (\texttt{message}, \texttt{products}, \texttt{draft\_\hspace{0pt}order},
\texttt{quick\_\hspace{0pt}replies}, \texttt{done}) qua \texttt{/api/\hspace{0pt}v1/\hspace{0pt}ai/\hspace{0pt}converse}.
Trên trình duyệt, \texttt{apps/web} cung cấp các trang dưới \texttt{tenants/\hspace{0pt}[tenantId]} và \texttt{/channels/\hspace{0pt}facebook/\hspace{0pt}callback}; \texttt{@hexta/sdk} quản lý access token ngắn
hạn (refresh token giữ hoàn toàn trong cookie \texttt{HttpOnly}, \texttt{SameSite=Lax}), bộ đọc SSE \texttt{readConsultEvents}, chuyển tiếp đơn nháp sang \texttt{checkout} và
giữ bản nháp cục bộ cùng khóa bất biến khi mạng gián đoạn.

\section{Kết quả và đánh giá}
\subsection{Tổng hợp kết quả kiểm thử và triển khai cô lập}
Các kết quả kiểm thử tự động, kiểm tra tĩnh, đồng bộ lược đồ và dựng cấu hình máy chủ đơn được tổng hợp trung thực từ báo cáo các phiên trong \texttt{agentic-memory} ở
Bảng~\ref{tab:eval-results}.

\begin{table}[htbp]
  \centering
  \footnotesize
  \setlength{\tabcolsep}{4pt}
  \renewcommand{\arraystretch}{1.2}
  \caption{Tổng hợp kết quả kiểm thử và xác minh trong môi trường cô lập}\label{tab:eval-results}
  \begin{tabularx}{\linewidth}{@{}
    >{\raggedright\arraybackslash}p{2.9cm}
    >{\raggedright\arraybackslash}X
    >{\raggedright\arraybackslash}p{2.1cm}@{}}
    \toprule
    Hạng mục & Lệnh thực thi, phạm vi và bằng chứng ghi nhận & Kết quả \\
    \midrule
    Go API, Shared và Connector (\texttt{-race}) &
    \texttt{go test -race -count=1 -p 1 ./services/api/... ./services/connector/... ./packages/shared/...} với PostgreSQL (\texttt{compose.connector-test.yml}) và Kafka cô lập: đạt toàn bộ \textbf{19 gói} ở đợt kiểm thử \textit{audit-hardening} (khi đó chạy trên Redpanda); sau khi chuyển sang Confluent \texttt{cp-\hspace{0pt}kafka:\hspace{0pt}7.6.1}, lần chạy lại ngày 18/11/2026 đạt \textbf{22/22 gói}, 203 ca \texttt{PASS}, 0 ca lỗi, chỉ bỏ qua ca nghiệm thu Telegram thật; các phiên tích hợp Facebook/Telegram đều đạt toàn bộ ca E2E và Outbox (gồm \texttt{TestConnectorFacebookToInbox}, \texttt{TestOutboxRelayPublishesCommittedEventToKafka}, \texttt{TestConnectorFacebookOAuthSendAndDeliveryResults}, \texttt{TestConnectorTelegramConnectIngressSendAndResults}). &
    PASS (19 gói) \\
    Dịch vụ AI Python &
    \texttt{./scripts/local.sh --profile test run --build --rm ai-tests} kết hợp \texttt{test-postgres} đã chạy migration Atlas: \textbf{109 passed, 0 skipped} ở đợt \textit{audit-hardening} (các đợt trước đạt 70/70 và 73/73 khi gắn DB; 70 passed, 6 skipped khi chưa gắn DSN). &
    PASS (109/109) \\
    Kiểm tra tĩnh, Atlas và OpenAPI &
    \texttt{sh scripts/check-gofmt.sh}, \texttt{go vet} (\texttt{api}/\texttt{connector}/\texttt{shared}/\texttt{test}), \texttt{make migrate-validate} (\texttt{api}/\texttt{ai}/\texttt{connector}), kiểm tra đồng bộ \texttt{atlas migrate diff} (gồm nâng cấp FK tổ hợp từ chối dữ liệu lệch), \texttt{make swagger} và kiểm thử hợp đồng router. &
    PASS \\
    Giao diện Web và bộ SDK &
    \texttt{pnpm --filter web lint}, \texttt{pnpm --filter web exec tsc --noEmit}, \texttt{pnpm --filter web build}; \texttt{pnpm --filter @hexta/sdk test} (\textbf{35 passed}), kiểm tra strict/no-\texttt{any} và build ESM/CJS/declaration; 7 kiểm thử \texttt{node --test} cho SSE và checkout. &
    PASS (SDK 35/35) \\
    Triển khai máy chủ đơn và GitOps &
    Dựng toàn bộ stack \texttt{compose.prod.yml} cục bộ với \texttt{DOMAIN=localhost} và khóa ngẫu nhiên (web/API/Swagger HTTP 200, webhook trả 403 khi thiếu chữ ký,
    \texttt{/internal/connections} trả 404 tại Caddy); kiểm thử \texttt{scripts/deploy-agent.sh} trong container Debian qua \textbf{9 kịch bản} cập nhật/rollback với Docker
    giả lập; \texttt{scripts/check-agent-instructions.py}. &
    PASS (9 kịch bản agent) \\
    \bottomrule
  \end{tabularx}
\end{table}

Ngoài kiểm thử tự động trên máy chủ giả lập, phiên kết nối Telegram (\texttt{2026-\hspace{0pt}11-\hspace{0pt}16\_\hspace{0pt}telegram-\hspace{0pt}bot-\hspace{0pt}connection}) còn kiểm chứng gọi trực tiếp Bot API thật tới bot thử nghiệm
\texttt{@hexta\_\hspace{0pt}test\_\hspace{0pt}bot} (token nạp từ tệp riêng tạm thời): \texttt{get\hspace{0pt}Me} trả về đúng định danh bot, \texttt{send\hspace{0pt}Message} tới chat không tồn tại trả về
\texttt{chat\_\hspace{0pt}not\_\hspace{0pt}found} (HTTP 400), token sai trả về \texttt{auth\_\hspace{0pt}expired} (HTTP 401) và không rò rỉ token trong thông điệp lỗi.

\subsection{Đánh giá đối chiếu theo các câu hỏi nghiên cứu}
\begin{itemize}
  \item \textbf{Đối với RQ1 (Quy trình tuần tự bất biến và hai chế độ bàn giao):} Các kiểm thử tích hợp trên PostgreSQL thật khẳng định việc tách
    \texttt{workflow\_\hspace{0pt}versions}/\hspace{0pt}\texttt{workflow\_\hspace{0pt}steps} khỏi \texttt{work\_\hspace{0pt}items} giữ nguyên cấu trúc bước của các công việc đang chạy khi sửa quy trình.
  Hai chế độ \texttt{direct} và \texttt{acceptance} xác định tường minh người chịu trách nhiệm ở từng bước và tuần tự hóa an toàn mọi tranh chấp
  \texttt{accept}/\texttt{assign} hay lặp \texttt{Idempotency-\hspace{0pt}Key}.
  \item \textbf{Đối với RQ2 (Định tuyến kịch bản/AI bám sát dữ liệu được cấp và kiểm soát hội thoại):} Bộ kiểm thử \texttt{services/ai} và \texttt{ConsultantService} xác nhận
    hệ thống chỉ chọn sản phẩm/facts từ ảnh chụp được Go cấp, loại bỏ mã giả mạo, tự động chuyển \texttt{human} khi thiếu dữ liệu hoặc chưa bật quyền sinh văn bản, và bảo vệ
    trạng thái hội thoại bằng \texttt{HandlingState()} cùng \texttt{control\_\hspace{0pt}generation}.
  Hàng rào \texttt{LLM\_\hspace{0pt}CALL\_\hspace{0pt}STARTED} ngăn gọi lặp mô hình trả phí khi phục hồi. Dù vậy, hiệu quả giảm thời gian chuẩn bị phản hồi và chất lượng câu chữ tiếng Việt từ mô hình
  thật chưa được đo trực tiếp trên dữ liệu và nhân viên thực tế.
  \item \textbf{Đối với RQ3 (Phân quyền, cô lập đa doanh nghiệp, kiểm soát cạnh tranh và lưu vết):} Tập kiểm thử có \texttt{-race} chứng minh mọi truy cập đều kiểm tra thành
    viên trực tiếp tại DB (\texttt{AuthorizeMember}, \texttt{ISessionRepository}), khóa ngoại tổ hợp \texttt{(tenant\_\hspace{0pt}id, id)}, phiên bản lạc quan (\texttt{version},
    \texttt{expected\_\hspace{0pt}version}, \texttt{context\_\hspace{0pt}revision}), khóa hàng giao dịch và ghi đồng bộ \texttt{audit\_\hspace{0pt}logs}.
  Việc cô lập ba tài khoản/CSDL và mã hóa AES-256-GCM AAD khóa chặt ranh giới bí mật giữa các dịch vụ.
\end{itemize}

\subsection{Ranh giới kiểm chứng và các hạng mục chưa thực thi}
Nhằm bảo đảm tính trung thực học thuật theo phạm vi phân định ở Chương~\ref{chap:introduction}, kết quả kiểm thử cô lập không được suy diễn thành nghiệm thu sản xuất. Các hạng
mục dưới đây hiện ở trạng thái \textbf{chưa thực thi (\texttt{NOT\_RUN})}:
\begin{enumerate}
  \item \textbf{Đo thời gian thao tác thực tế của nhân viên và gọi LLM thật:} Các khóa mô hình thật (\texttt{OPENAI\_\hspace{0pt}API\_\hspace{0pt}KEY}, \texttt{GEMINI\_\hspace{0pt}API\_\hspace{0pt}KEY},
    \texttt{ANTHROPIC\_\hspace{0pt}API\_\hspace{0pt}KEY}) để trống trong quá trình kiểm thử;
  lệnh \texttt{smoke.py --live}, đánh giá chất lượng câu trả lời từ mô hình thương mại, và thực nghiệm đo thời gian chuẩn bị phản hồi của nhân viên trên tập dữ liệu 10 cây
  (đối chứng với trả lời nhanh) đều là \texttt{NOT\_RUN}.
  \item \textbf{Nghiệm thu webhook công khai từ Meta App Review và Telegram trực tiếp:} Luồng nhận webhook HTTPS công khai từ Meta/Telegram thật và kiểm thử cửa sổ 24 giờ trên
    trang thật là \texttt{NOT\_RUN} do phụ thuộc quy trình xét duyệt Meta App Review và tên miền công khai đang vận hành.
  \item \textbf{Triển khai VPS thật, kiểm thử tải và bộ test liên dịch vụ \texttt{test/...}:} Cấu hình \texttt{compose.\hspace{0pt}prod.\hspace{0pt}yml} mới chạy thử với \texttt{DOMAIN=\hspace{0pt}localhost},
    còn \texttt{deploy-\hspace{0pt}agent.\hspace{0pt}sh} mới kiểm thử trong container giả lập Docker;
  vận hành trên VPS 2 vCPU / 16 GB RAM thật với chứng chỉ Let's Encrypt thật, diễn tập khôi phục sao lưu, bộ kiểm thử \texttt{go test ./\hspace{0pt}test/...} (cần cụm máy chủ chạy đồng
  thời) và kiểm thử tải là \texttt{NOT\_RUN}. Bộ mục tiêu chịu tải T1--T5 (Bảng~\ref{tab:load-targets}) đã được xác lập cho cấu hình hai máy (máy hạ tầng 2 vCPU / 4 GB và máy ứng dụng 2 vCPU / 16 GB) nhưng chưa có số đo trên máy chủ thật; mốc xuất phát khoảng 1 tin/giây mới chỉ được đo cục bộ ở bước ghi Kafka; việc tách triển khai thành hai máy cũng chưa được cài đặt.
  \item \textbf{Duyệt giao diện thủ công toàn trình trên trình duyệt:} Ngoài kiểm thử khói bán hàng ban đầu và kiểm tra lint/build tự động, việc duyệt thủ công đa tài khoản
    trên trình duyệt đối với quy trình công việc nhiều bước, hộp thư Desk, cấu hình kênh Facebook/Telegram, tương thích di động và thao tác bàn phím vẫn là \texttt{NOT\_RUN}.
\end{enumerate}
